% GENERATED FILE. Edit canonical lessons, then run npm run generate:books.
% canonical-source-hash: fnv1a64:aaa39820831c4cfb
% canonical-lessons: GU-C380-purtu, GU-C380-tangi, GU-C380-thodu, GU-C380-sahej, GU-C380-pahoci-valvu

\chapter{Enough or Not}
\label{ch:gu-enough-or-not}
\hlchaptermodality{380}

\begin{chapteropening}
\textbf{By the end of this chapter:} \emph{I can say whether there will be enough.}

Complete the last lesson of chapter 380: I can say whether there will be enough.
\end{chapteropening}

\section[\texorpdfstring{\gu{પૂરતું} (pūrtũ)}{પૂરતું (pūrtũ)}]{\gu{પૂરતું} (pūrtũ) — enough}

\label{lesson:GU-C380-purtu}



\begin{quote}
Before the new one: say the Gujarati for absolutely, then the Gujarati for solely.
\end{quote}

\subsection*{You'll want to know: \gu{પૂરતું}}
\textbf{\gu{પૂરતું}} — \emph{pūrtũ} — \textquotedblleft{}enough\textquotedblright{}.

\begin{quote}
\gu{જો} \gu{પૂરતો} \gu{વરસાદ} \gu{પડશે} \gu{તો} \gu{પાણીની} \gu{ચિંતા} \gu{નહીં} \gu{રહે}.
\end{quote}

\begin{quote}
\emph{jo pūrto varsād paḍśe to pāṇīnī cintā nahī̃ rahe.}
\end{quote}

\begin{quote}
If there is enough rain, there will be no worry about water.
\end{quote}

\begin{cousinweb}
\textbf{\gu{પૂરતું}} (\emph{pūrtũ}) is built on \emph{pūrũ}, \textquotedblleft{}full\textquotedblright{}, which CDIAL traces to Sanskrit \emph{pūra}, \textquotedblleft{}filling\textquotedblright{} (CDIAL 8330): as much as fills the need.
\end{cousinweb}

\begin{grammarlens}[title={What you've built}]
One more word for saying what happens if.
\end{grammarlens}

\subsection*{Your turn}
\begin{itemize}
\raggedright
  \item \emph{Say it:} \emph{pūrtũ}
  \item \emph{Say it:} \emph{pūrtũ}, once more
  \item \emph{Say it:} \emph{jo pūrto varsād paḍśe to pāṇīnī cintā nahī̃ rahe.} — If there is enough rain, there will be no worry about water.
  \item \emph{Recall:} say the Gujarati for absolutely, then the Gujarati for solely, then say \emph{pūrtũ} again
\end{itemize}

\begin{tcolorbox}[breakable,colback=teal!4,colframe=teal!35!black,title={Before you move on}]
What does \gu{પૂરતું} mean? (\textquotedblleft{}Enough\textquotedblright{}.) Say it once more.
\end{tcolorbox}
\section[\texorpdfstring{\gu{તંગી} (tangī)}{તંગી (tangī)}]{\gu{તંગી} (tangī) — a shortage}

\label{lesson:GU-C380-tangi}



\begin{quote}
Before the new one: say the Gujarati for solely, then the Gujarati for enough.
\end{quote}

\subsection*{You'll want to know: \gu{તંગી}}
\textbf{\gu{તંગી}} — \emph{tangī} — \textquotedblleft{}a shortage\textquotedblright{}.

\begin{quote}
\gu{જો} \gu{ઉનાળામાં} \gu{વરસાદ} \gu{મોડો} \gu{આવશે} \gu{તો} \gu{પાણીની} \gu{તંગી} \gu{પડશે}.
\end{quote}

\begin{quote}
\emph{jo unāḷāmā̃ varsād moḍo āvśe to pāṇīnī tangī paḍśe.}
\end{quote}

\begin{quote}
If the rain comes late this summer, there will be a water shortage.
\end{quote}

A feminine noun.

\begin{cousinweb}
\textbf{\gu{તંગી}} (\emph{tangī}) is usually traced to Persian \emph{tangī}, \textquotedblleft{}narrowness, want\textquotedblright{}.
\end{cousinweb}

\begin{grammarlens}[title={What you've built}]
One more word for saying what happens if.
\end{grammarlens}

\subsection*{Your turn}
\begin{itemize}
\raggedright
  \item \emph{Say it:} \emph{tangī}
  \item \emph{Say it:} \emph{tangī}, once more
  \item \emph{Say it:} \emph{jo unāḷāmā̃ varsād moḍo āvśe to pāṇīnī tangī paḍśe.} — If the rain comes late this summer, there will be a water shortage.
  \item \emph{Recall:} say the Gujarati for solely, then the Gujarati for enough, then say \emph{tangī} again
\end{itemize}

\begin{tcolorbox}[breakable,colback=teal!4,colframe=teal!35!black,title={Before you move on}]
What does \gu{તંગી} mean? (\textquotedblleft{}A shortage\textquotedblright{}.) Say it once more.
\end{tcolorbox}
\section[\texorpdfstring{\gu{થોડું} (thoḍũ)}{થોડું (thoḍũ)}]{\gu{થોડું} (thoḍũ) — a little}

\label{lesson:GU-C380-thodu}



\begin{quote}
Before the new one: say the Gujarati for enough, then the Gujarati for a shortage.
\end{quote}

\subsection*{You'll want to know: \gu{થોડું}}
\textbf{\gu{થોડું}} — \emph{thoḍũ} — \textquotedblleft{}a little\textquotedblright{}.

\begin{quote}
\gu{જો} \gu{તમે} \gu{થોડી} \gu{રાહ} \gu{જોશો} \gu{તો} \gu{બસ} \gu{આવશે}.
\end{quote}

\begin{quote}
\emph{jo tame thoḍī rāh jośo to bas āvśe.}
\end{quote}

\begin{quote}
If you wait a little, the bus will come.
\end{quote}

\begin{cousinweb}
\textbf{\gu{થોડું}} (\emph{thoḍũ}) is inherited from Sanskrit \emph{stōka}, \textquotedblleft{}a drop\textquotedblright{}, and as an adverb \textquotedblleft{}a little\textquotedblright{}; CDIAL lists the Gujarati word, \textquotedblleft{}a little, few\textquotedblright{} (CDIAL 13720).
\end{cousinweb}

\begin{grammarlens}[title={What you've built}]
One more word for saying what happens if.
\end{grammarlens}

\subsection*{Your turn}
\begin{itemize}
\raggedright
  \item \emph{Say it:} \emph{thoḍũ}
  \item \emph{Say it:} \emph{thoḍũ}, once more
  \item \emph{Say it:} \emph{jo tame thoḍī rāh jośo to bas āvśe.} — If you wait a little, the bus will come.
  \item \emph{Recall:} say the Gujarati for enough, then the Gujarati for a shortage, then say \emph{thoḍũ} again
\end{itemize}

\begin{tcolorbox}[breakable,colback=teal!4,colframe=teal!35!black,title={Before you move on}]
What does \gu{થોડું} mean? (\textquotedblleft{}A little\textquotedblright{}.) Say it once more.
\end{tcolorbox}
\section[\texorpdfstring{\gu{સહેજ} (sahej)}{સહેજ (sahej)}]{\gu{સહેજ} (sahej) — slightly}

\label{lesson:GU-C380-sahej}



\begin{quote}
Before the new one: say the Gujarati for a shortage, then the Gujarati for a little.
\end{quote}

\subsection*{You'll want to know: \gu{સહેજ}}
\textbf{\gu{સહેજ}} — \emph{sahej} — \textquotedblleft{}slightly\textquotedblright{}.

\begin{quote}
\gu{જો} \gu{તમે} \gu{સહેજ} \gu{ખસશો} \gu{તો} \gu{અમે} \gu{પણ} \gu{બેસી} \gu{શકીશું}.
\end{quote}

\begin{quote}
\emph{jo tame sahej khasśo to ame paṇ besī śakīśũ.}
\end{quote}

\begin{quote}
If you move up slightly, we will be able to sit too.
\end{quote}

\begin{cousinweb}
\textbf{\gu{સહેજ}} (\emph{sahej}) is usually traced to Sanskrit \emph{sahaja}, \textquotedblleft{}born together, innate, natural\textquotedblright{} (Monier-Williams); \textquotedblleft{}slightly, easily\textquotedblright{} is a Gujarati sense the source does not give.
\end{cousinweb}

\begin{grammarlens}[title={What you've built}]
One more word for saying what happens if.
\end{grammarlens}

\subsection*{Your turn}
\begin{itemize}
\raggedright
  \item \emph{Say it:} \emph{sahej}
  \item \emph{Say it:} \emph{sahej}, once more
  \item \emph{Say it:} \emph{jo tame sahej khasśo to ame paṇ besī śakīśũ.} — If you move up slightly, we will be able to sit too.
  \item \emph{Recall:} say the Gujarati for a shortage, then the Gujarati for a little, then say \emph{sahej} again
\end{itemize}

\begin{tcolorbox}[breakable,colback=teal!4,colframe=teal!35!black,title={Before you move on}]
What does \gu{સહેજ} mean? (\textquotedblleft{}Slightly\textquotedblright{}.) Say it once more.
\end{tcolorbox}
\section[\texorpdfstring{\gu{પહોંચી} \gu{વળવું} (pahõcī vaḷvũ)}{પહોંચી વળવું (pahõcī vaḷvũ)}]{\gu{પહોંચી} \gu{વળવું} (pahõcī vaḷvũ) — to cope}

\label{lesson:GU-C380-pahoci-valvu}



\begin{quote}
Before the new one: say the Gujarati for a little, then the Gujarati for slightly.
\end{quote}

\subsection*{You'll want to know: \gu{પહોંચી} \gu{વળવું}}
\textbf{\gu{પહોંચી} \gu{વળવું}} — \emph{pahõcī vaḷvũ} — \textquotedblleft{}to cope\textquotedblright{}.

\begin{quote}
\gu{જો} \gu{બધા} \gu{મદદ} \gu{કરશે} \gu{તો} \gu{આપણે} \gu{આ} \gu{કામને} \gu{પહોંચી} \gu{વળીશું}.
\end{quote}

\begin{quote}
\emph{jo badhā madad karśe to āpaṇe ā kāmne pahõcī vaḷīśũ.}
\end{quote}

\begin{quote}
If everyone helps, we will cope with this work.
\end{quote}

Past (he …): \gu{તે} \gu{બધી} \gu{મુશ્કેલીઓને} \gu{પહોંચી} \gu{વળ્યો}, \emph{te badhī muśkelīone pahõcī vaḷyo}.

\begin{cousinweb}
\textbf{\gu{પહોંચી} \gu{વળવું}} (\emph{pahõcī vaḷvũ}) is built from \textbf{\gu{પહોંચવું}} (\emph{pahõcvũ}), \textquotedblleft{}to reach\textquotedblright{} (CDIAL 8716), and \textbf{\gu{વળવું}} (\emph{vaḷvũ}), \textquotedblleft{}to turn\textquotedblright{} (CDIAL 11405): to reach a task and turn it round.
\end{cousinweb}

\begin{grammarlens}[title={What you've built}]
One more word for saying what happens if.
\end{grammarlens}

\subsection*{Your turn}
\begin{itemize}
\raggedright
  \item \emph{Say it:} \emph{pahõcī vaḷvũ}
  \item \emph{Say it:} \emph{pahõcī vaḷvũ}, once more
  \item \emph{Say it:} \emph{jo badhā madad karśe to āpaṇe ā kāmne pahõcī vaḷīśũ.} — If everyone helps, we will cope with this work.
  \item \emph{Recall:} say the Gujarati for a little, then the Gujarati for slightly, then say \emph{pahõcī vaḷvũ} again
\end{itemize}

\begin{tcolorbox}[breakable,colback=teal!4,colframe=teal!35!black,title={Before you move on}]
What does \gu{પહોંચી} \gu{વળવું} mean? (\textquotedblleft{}To cope\textquotedblright{}.) Say it once more.
\end{tcolorbox}
